\chapter*{Glossary} \phantomsection \addcontentsline{toc}{chapter}{Glossary} \markboth{Glossary}{Glossary}

{\small
\setlength{\LTpre}{0pt} \setlength{\LTpost}{0pt} \setlength{\tabcolsep}{4pt} \renewcommand{\arraystretch}{1.1} \begin{longtable}{@{}>{\raggedright\arraybackslash\bfseries} p{0.27\textwidth}>{\raggedright\arraybackslash} p{0.69\textwidth}@{}} Active fund, Active management & A collective fund where securities are bought and sold actively to try and outperform the general market. Normally more expensive than passive funds, as active fund managers think they can generate, and charge for, alpha. Hedge funds are an extreme example of active funds. As distinct from passive funds and passive management. See page 106. \\[0.65em]
Alpha & The returns of an active manager or trader can be split into beta and alpha. The beta are the returns you could get from investing in the general market, i.e. in a passive index fund. Any additional return due to the manager’s skill is alpha. \\[0.65em]
Alternative beta & A kind of beta, but which requires active trading to achieve. So for example to earn the equity value premium, which is the return from being long low price:earnings (PE) and short high PE equities, you need to buy and sell the appropriate shares at the right time. Like beta, and unlike alpha, this kind of return does not require skill. \\[0.65em]
Annualised cash volatility target & See volatility target. \\[0.65em]
Asset allocating investor & An investor who usually does not forecast asset prices, but uses a systematic framework to create the best possible portfolio. See page ix and page 225. \\[0.65em]
Asset class & Name for a general kind of investable security, e.g. equities, bonds, commodity futures and commercial property are all asset classes. \\[0.65em]
Back-testing & A back-test is a historical simulation of what performance and behaviour would have been for a trading system, based on the data and prices that occurred in the past. See page 14. \\[0.65em]
Beta & Used as a shorthand for getting returns from being exposed to the market generally, without having any particular skill; as distinct from alpha. Owning a passive index fund, such as an ETF linked to the French CAC 30 stock market, will give you beta returns for the relevant market. See also alternative beta. Also: to what extent a particular asset is exposed to the overall market. Low beta assets have less exposure to the overall market, and are safer, than high beta assets. \\[0.65em]
Block value & The increase in value of one instrument block when its price goes up by one percentage point. See page 154. \\[0.65em]
Bootstrapping & A method of portfolio optimisation which takes into account the uncertainty in historical data, unlike single period optimisation. See page 75. \\[0.65em]
Calibration & Fine tuning a trading rule by looking at the performance or behaviour of variations during back-testing; a kind of fitting. See page 52. \\[0.65em]
Carry & A trading strategy where you profit from the difference between yields and funding costs if prices remain unchanged. A recommended carry trading rule is discussed in chapter seven, on page 119 and in appendix B, from page 285. See also FX Carry. \\[0.65em]
Cash volatility target & See volatility target. \\[0.65em]
Cognitive bias & A psychological flaw in human thought processes which results in poor decisions being made. See page 12. \\[0.65em]
Collective fund & A kind of fund which buys you a share in a portfolio of securities. These can be exchange traded funds, mutual funds in the US, or in the UK investment trusts and unit trusts. Collective funds can be active or passive. See page 106. \\[0.65em]
Combined forecast & The forecast you get for a particular instrument after combining the individual forecasts from multiple trading rules and multiplying by the forecast diversification multiplier. See page 125. \\[0.65em]
Correlation, Correlated & A measure of how two things co-move. Normally used in relation to daily returns from assets or profits from trading rules. A correlation of -1 indicates two things always move in opposite directions, +1 indicates they always move in the same direction and 0 means there is no linear relationship (uncorrelated). \\[0.65em]
Daily cash volatility target & See volatility target. \\[0.65em]
Data first & A kind of fitting where you specify a very general trading rule or model of price movement with many parameters and then use a statistical method to find the best choice of parameters from the available data. Opposite to ideas first. See page 26. \\[0.65em]
Data mining & The process of fitting to historical data in such a way that you will always find at least one apparently profitable trading rule. For example, if you run thousands of back-tests for different rules over a particular historical period, at least one will look great. Because they are highly tuned to fit the past, data mined trading rules usually perform badly in a future that is never the same. Also see over-fitting. \\[0.65em]
Derivative & A way of benefiting from an asset going up or down in price without actually owning it. Futures, spread bets, options and contracts for difference are all examples of derivatives. Derivatives often provide more leverage than the underlying assets. \\[0.65em]
Diversification multiplier & A portfolio will have a lower standard deviation of returns than its individual assets, assuming they are not perfectly correlated. This assumes the assets all have the same standard deviation of returns. The diversification multiplier multiplies portfolio returns to get them back to the same standard deviation as the underlying assets. It is always 1, for portfolios where all assets are perfectly correlated, or larger than 1. See instrument diversification multiplier and forecast diversification multiplier. See page 129. \\[0.65em]
Dynamic & A strategy where we actively buy and sell assets to express opinions on risk adjusted asset returns. As opposed to static. See page 38. \\[0.65em]
Equity value & A type of relative value trading in the equity market where you buy ‘cheap’ stocks with low price earnings (PE), high dividend yield or similar value characteristics, and sell ‘expensive’ stocks with high PE, low yield or similar. \\[0.65em]
Equity volatility index & A measure of how volatile an equity market is expected to be, derived from option prices. Sometimes called the ‘fear index’. The key European index is the V2TX and in the US it’s the VIX. Futures can be traded on each of these indices. See skew concept box, page 34. \\[0.65em]
Exchange traded fund (ETF) & Usually a type of passive indexed, collective fund which can be traded and owned like a normal equity, and has low fees. Some ETFs implement dynamic strategies and are more expensive. \\[0.65em]
Execution cost & Part of the cost of trading an instrument, equal to the difference between the mid-price and the traded price. Those who do not trade in large sizes can assume they will pay the bid (if selling) or offer (if buying), which implies it will cost them half the difference between the bid and the offer to trade. See page 179. \\[0.65em]
Expanding window & A way of doing out of sample testing. You pretend you are at a point in the past and then use data from the beginning of your price history up to that point to choose trading rules or portfolio weights. You then test those rules or weights on the data for a year after that point. Then you move forward a year to the next point and repeat. See page 56. \\[0.65em]
Exponentially weighted moving average (EWMA) & A kind of moving average where you weight more recent observations more heavily. If \(P_t\) is the most recent data point then the EWMA is \((A \times P_t) + (A \times (1 - A) \times P_{t-1}) + (A \times (1 - A)^2 \times P_{t-2}) + \dots\). You can also calculate this recursively. If you have yesterday’s EWMA \(E_{t-1}\) then today’s EWMA is \((A \times P_t) + (E_{t-1} \times (1 - A))\). Used in the EWMAC trading rule and also in the estimation of price volatility. \\[0.65em]
Exponentially weighted moving average crossover (EWMAC) & A trend following trading rule where you take a fast exponentially weighted moving average (EWMA) of the price with a short look-back, and a slow EWMA with a longer look-back. If the fast EWMA is above the slow then prices have been rising and you buy, and vice versa. See page 117 and page 282 in appendix B. \\[0.65em]
Fitting & The process of picking the best trading rule. Usually this would involve running back-tests over historical data and finding the most profitable rule. See page 52. See also ideas first and data first, and over-fitted. \\[0.65em]
Forecast & An estimate of how much an instrument’s price will go up or down, translated into a numeric scale. The forecast can either be the discretionary forecast of a semi-automatic trader or the fixed forecast of the asset allocating investor; for staunch systems traders a single systematic trading rule or a combined forecast blending different trading rules. See page 110. \\[0.65em]
Forecast diversification multiplier & The diversification multiplier required so that the combined forecast for a particular instrument has the expected average absolute value (recommended value: 10). See page 129. \\[0.65em]
Forecast scalar & A value you multiply a forecast by to ensure it has the right average absolute value (recommended value: 10). See page 297. \\[0.65em]
Forecast weights & The weights used to combine forecasts from multiple trading rules into a single combined forecast for a particular instrument. See page 126. \\[0.65em]
Foreign exchange (FX) carry & A foreign exchange (FX) trading strategy where you borrow in low interest currencies, exchange the money for a high interest currency, put it into the bank, and earn the difference. This strategy will fail if the high interest currency depreciates by more than the gap in interest rates. See also carry. \\[0.65em]
Framework & The process I describe in this book to translate forecasts into actual investment decisions. The framework includes components for combining trading rules and instruments, volatility targeting and position sizing. Together with any trading rules the framework makes up the trading system. See page 96. \\[0.65em]
Fundamental data & Non-price data used for forecasting, e.g. earnings of a firm, inflation rates or weather forecasts. Opposite of technical data. See page 43. \\[0.65em]
Gaussian normal distribution & A statistical distribution which is shaped like a bell, often assumed for returns. If your returns are Gaussian normal then you will see returns one sigma or less around the average about 68\% of the time, and returns two sigma or less about 95\% of the time. See page 21. \\[0.65em]
Half-Kelly & See Kelly criterion. \\[0.65em]
Handcrafted optimisation & A method of portfolio optimisation where you set weights by hand, grouping assets together and using only estimates of correlations (and optionally, Sharpe ratios). See page 78. \\[0.65em]
Ideas first & A type of fitting where you conceive an idea, choose a specific trading rule and then back-test the rule on historical data. Opposite to data first. See page 26. \\[0.65em]
In sample & Using historical data to select trading rules or portfolio weights and then testing those rules on the same data, which means including future information. A bad thing. Opposite to out of sample. See page 54. \\[0.65em]
Index tracker & See passive fund. \\[0.65em]
Instrument & Something you trade or invest in, e.g. an equity like Apple stock, an oil futures contract or a spread bet on the EUR/USD FX rate. See page 101. \\[0.65em]
Instrument block & ‘One’ of an instrument – the minimum discrete unit you can economically trade an instrument in, which might be 100 shares, one futures contract or £1 a point on a spread bet. See page 154. \\[0.65em]
Instrument currency volatility & The expected daily standard deviation of returns of owning one instrument block, measured in the currency of the instrument. Equal to price volatility multiplied by block value. See page 158. \\[0.65em]
Instrument diversification multiplier & A type of diversification multiplier that accounts for the diversification across the returns of trading subsystems. See page 169. \\[0.65em]
Instrument value volatility & The expected daily standard deviation of returns of owning one instrument block, measured in the same currency as the trader or investor’s trading capital. Equal to instrument currency volatility multiplied by the relevant exchange rate. See page 158. \\[0.65em]
Instrument weights & In my framework the weights that different trading subsystems, each trading a single instrument, have in a portfolio of subsystems. See page 165. \\[0.65em]
Kelly criterion & A way of determining the percentage volatility target you should use, given your expectation of Sharpe ratio (SR). A SR of 0.5 implies you should use a risk percentage of 50\%. However it is safer to use Half-Kelly and set the risk to half what Kelly requires, i.e. for an SR of 0.5 set risk to 25\%. See page 143. \\[0.65em]
Law of active management & An equation which states that the information ratio of a trading strategy is proportional to the square root of the number of independent bets made each year. See page 42. \\[0.65em]
Leverage & Borrowing to invest, either in an explicit way or by using a derivative such as a future or spread bet where your exposure is greater than your initial cash payment. \\[0.65em]
Liquidity & How easy it is to trade quickly in size without changing the price significantly. Shares in blue chip companies are liquid because you can usually sell \$1 million worth within minutes at close to the prevailing price. A \$1 million house cannot be sold in a single day without giving buyers a substantial discount from what it could achieve given more time. The blue chip shares are liquid; the house is not. See page 35. \\[0.65em]
Mean reversion & Any trading strategy where you assume asset prices will revert to an equilibrium or fair value. For example, you might think 1.70 is a fair value for the price of the GBP/USD FX rate. If the rate goes below 1.70 you’ll buy pounds and sell dollars, and if it goes higher you’d sell pounds (and buy dollars). See also relative value. \\[0.65em]
Merger arbitrage & A trading strategy, where if a merger or takeover is occurring you buy the company to be acquired, and usually short sell the acquirer as a hedge. Normally this is done at a discount to the takeover price, the gap reflecting uncertainty about whether a deal will go ahead. Profits can be made by selecting deals where the gap is too large given the level of uncertainty. This is a negative skew style of trading, because occasional large losses are made when deals fall through. \\[0.65em]
Momentum & When asset prices go up after previous rises; or go down after falls. Trend following rules try and profit from momentum. \\[0.65em]
Moving average & A weighted average of a value over the last N points in time, where N is the look-back window. The weighting can either be a simple average, with all points equally weighted, or an exponentially weighted moving average, with more recent points given a higher weight. Used in the estimation of price volatility. \\[0.65em]
Open interest & The number of futures contracts outstanding for a particular instrument. A measure of how actively an instrument is traded. \\[0.65em]
Out of sample & Using historical data to select trading rules or portfolio weights without looking into the future. A good thing, in contrast to in sample. See page 54. \\[0.65em]
Over-fitted & When something is fitted in such a way that it matches the past data too well and is unlikely to perform well when actually traded. Sometimes referred to as curve fitting. Also see data mining. Refer to page 52. \\[0.65em]
Passive fund, Passive index, Passive management & A type of collective fund which is invested and managed passively in a portfolio of securities according to some predefined weights. Usually these funds are index trackers which follow indices such as the FTSE 100 or S\&P 500. Passive funds are a cheap way to get broad market, or beta, exposure. As distinct from active funds. See page 106. \\[0.65em]
Percentage volatility target & The target expected annualised percentage risk your trading system is exposed to. When multiplied by trading capital we get the annualised cash volatility target. See also volatility target and page 138. \\[0.65em]
Portfolio optimisation & The process of finding the optimal (best) portfolio by deciding in which portfolio weights to hold your assets. Normally a portfolio consists of positions in various assets, e.g. stocks and bonds. In my framework portfolios can be of (a) multiple trading rules for which we find forecast weights, and for (b) trading subsystems, one per instrument, for which we optimise instrument weights. You should use handcrafting or bootstrapping to calculate portfolio weights; I do not recommend single period optimisation. See page 70. \\[0.65em]
Portfolio weighted position & The position you hold in an instrument at the trading system level. Because a trading system consists of a portfolio of trading subsystems, one for each instrument, this is equal to the instrument’s subsystem position multiplied by the instrument weight and the instrument diversification multiplier. See page 173. \\[0.65em]
Portfolio weights & In general the weights that assets have in a portfolio. See instrument weights and forecast weights. \\[0.65em]
Position inertia & If the current position is less than 10\% away from the rounded target position then to reduce costs you shouldn’t trade. See page 174. \\[0.65em]
Predictable risk & The component of volatility in asset prices or portfolio returns which you can predict using estimates of standard deviation and correlation. See unpredictable risk and page 39. \\[0.65em]
Price volatility & The expected standard deviation of instrument price daily returns, in percentage points per day. See page 155. \\[0.65em]
Relative value & Any trading strategy where you buy something that is cheap, and sell something similar that is expensive as a hedge. For example, see equity value. Often a negative skew strategy, as we suffer occasional large losses. \\[0.65em]
Risk parity & A type of passive fund where assets are held in equal shares according to expected risk. See page 38. \\[0.65em]
Risk premium & The reward you get from investing in something that carries some risk. For example, equities are riskier than bonds so we should expect higher returns to compensate us. See page 31. \\[0.65em]
Rolling window & A method of out of sample testing. You pretend you are at a point in the past and then use data from the last N years (where N can be any value) up to that point to choose trading rules or portfolio weights. You then test those rules or weights on data for a year after that point. Then you move forward a year to the next point and repeat. See page 56. \\[0.65em]
Semi-Automatic Trader & A trader who makes their own forecasts about price movements in a discretionary fashion, but then incorporates them into a systematic framework. See page ix and page 209. \\[0.65em]
Sharpe ratio (SR) & A measure of how profitable a trading strategy is, with returns adjusted for risk. Formally it is the mean return over some time period divided by the standard deviation of returns over the same time period. In this book I normally use the annualised Sharpe ratio – annualised returns divided by annualised standard deviation. If we’re not using derivatives the ‘risk free’ interest rate should be deducted from annualised returns before calculating the Sharpe ratio. See page 32. \\[0.65em]
Short selling & A way of benefiting from an asset’s fall in price, usually used for equity trading. You would usually borrow the stock with a loan agreement before selling it. On closing the trade the stock is returned and you’ll profit if the price has fallen since inception. \\[0.65em]
Sigma & Shorthand for one unit of standard deviation. See page 21. \\[0.65em]
Single period optimisation & The classic way of performing a portfolio optimisation: using a single average over historical data of asset return means, standard deviations and correlations; these estimates are then used to do a single optimisation. \\[0.65em]
Skew & A measure of how symmetric the returns from an asset or trading rule are. Positive skew means you get more losses and fewer gains. But the average loss is smaller than the average gain. Negative skew means you get more gains and fewer losses, and average losses are larger than average gains. See page 32. \\[0.65em]
Speed limit & A limit on how quickly you should trade an instrument. The speed limit will depend on the standardised cost of the instrument, and how much you are prepared to pay in Sharpe ratio (SR) units per year in costs. I recommend paying no more than 0.13 SR units if you are a staunch systems trader and 0.08 SR units for semi-automatic traders and asset-allocating investors. See page 187. \\[0.65em]
Standard deviation & A measure of how dispersed some data is around its average value. If your data points are \(x_1, x_2, \ldots, x_n\) then the average \(x^* = (x_1 + x_2 + \ldots + x_n) \div n\). The standard deviation is \(\sqrt{(1 \div n)[(x_1 - x^*)^2 + (x_2 - x^*)^2 + \ldots + (x_n - x^*)^2]}\). Often applied to daily returns in prices or trading system profits. An approximate way to calculate an annualised standard deviation of returns is to multiply the daily standard deviation by 16, the square root of 256, approximately the number of trading days in one year. See page 21. \\[0.65em]
Standardised cost & A volatility standardised method of measuring costs which is comparable across instruments. To calculate, add up all the costs of trading \(C\), including execution costs, fees and taxes. Then double to get the cost of a round trip, a buy and a sell, and divide by the annualised instrument cash volatility \(ICV\) of the instrument. The formula is \(2 \times C \div (16 \times ICV)\). See page 181. \\[0.65em]
Static & A strategy where you do not actively trade the assets you are investing in because of changes in expected risk adjusted returns. As opposed to dynamic. See page 38. \\[0.65em]
Staunch Systems Trader & A trader who uses one or more systematic trading rules and incorporates them into a systematic framework. See page x and page 245. \\[0.65em]
Subsystem position & The position nominally held by a single instrument’s trading subsystem, trading only one instrument with the entire trading capital, given a forecast of returns. Equal to the volatility scalar multiplied by forecast, divided by 10. See page 159. \\[0.65em]
Survivorship bias & A problem where assets that disappeared are missing from historic data sets, such as shares in firms which went bankrupt. Because we don’t see the losses from these instruments we are likely to overestimate how profitable investing really was in the past. \\[0.65em]
Technical & A kind of trading rule that uses only price data to predict prices, such as trend following; no fundamental data is used. See page 43. \\[0.65em]
Trading capital & The amount of capital at risk in your trading system. Also see volatility target. See page 138. \\[0.65em]
Trading rules & A systematic rule used to predict whether the price of an instrument will go up or down, and by how much. When one or more trading rules are put into a framework they form a complete trading system for trading systematically. See page 109. \\[0.65em]
Trading subsystem & A trading subsystem is a notional part of a larger trading system which trades a single instrument. See page 103. \\[0.65em]
Trading system & A set of one or more trading rules within a framework creates a complete trading system; it makes forecasts about instrument price movements and translates these into positions and hence trades suitable for a particular level of trading capital. For asset allocating investors and staunch systems traders the trading system is made up of a portfolio of trading subsystems, one per instrument, for a fixed set of instruments and trading rules (although asset allocating investors use only one rule with a constant forecast). For semi-automatic traders the system is made up of discretionary forecasts for an ad hoc group of instruments, so the make-up of the group of trading subsystems will vary over time. See page 98. \\[0.65em]
Trend following & A technical trading strategy which tries to capture momentum in price; you buy things going up and sell things that have fallen in price. Tends to have positive skew. \\[0.65em]
Turnover & A way of measuring trading speed. Turnover is measured in round trips per year, where a round trip consists of a buy and a sell of an average sized position. So a trading system with a turnover of 10 units is expected to do ten buys and ten sells in a year. See page 184. \\[0.65em]
Unpredictable risk & The component of volatility in asset prices or portfolio returns which you didn’t or couldn’t predict in advance. This could be because standard deviations or correlations changed, or because your risk model was wrong. See predictable risk and page 39. \\[0.65em]
Variation & A variation on a trading rule has different parameter value(s) but is otherwise the same. For example ‘buy in the range 1.5 to 2.0’ is a variation of the general rule of ‘buy in the range X to Y’. See page 53. \\[0.65em]
Volatility & A shorthand term for standard deviation. \\[0.65em]
Volatility look-back period & Period of recent time used to calculate a standard deviation of returns. Used to calculate the price volatility and in trading rules like EWMAC. \\[0.65em]
Volatility scalar & A parameter which accounts for the difference in the risk you want for your portfolio (volatility target) relative to the risk of an instrument. It’s the number of instrument blocks you would hold if you invested your entire trading capital into one instrument, with a forecast of +10. Equal to daily cash volatility target divided by instrument value volatility. See page 159. \\[0.65em]
Volatility standardisation & Adjusting returns or costs so they have the same expected risk. See page 40. \\[0.65em]
Volatility target & The target expected standard deviation of returns for the entire trading system. Can be expressed in cash terms, in the same currency as trading capital, either as the annualised cash volatility target or daily cash volatility target. The daily cash volatility target is the annualised cash volatility target divided by ‘the square root of time’, which assuming a 256 business day year is 16. It can also be expressed as a percentage volatility target, as a percentage of trading capital. The cash targets are equal to trading capital multiplied by the percentage volatility target. See page 137. \\[0.65em]
\end{longtable}
}
